\documentclass[11pt,a4paper]{article}

% -----------------------------------------------------------------------------
% Paquetes: todos forman parte de una instalacion estandar de LaTeX.
% -----------------------------------------------------------------------------
\usepackage[utf8]{inputenc}
\usepackage{enumitem}
\usepackage{amsmath}
\usepackage{amssymb}
\usepackage{booktabs}
\usepackage[T1]{fontenc}
% Fecha en espanol sin depender de babel ni de spanish.ldf.
\newcommand{\FechaEnEspanol}{%
  \number\day\space de\space
  \ifcase\month
    \or enero\or febrero\or marzo\or abril\or mayo\or junio\or
    julio\or agosto\or septiembre\or octubre\or noviembre\or diciembre
  \fi\space de\space\number\year
}
\usepackage[a4paper,margin=2.5cm,headheight=42pt]{geometry}
\usepackage{graphicx}
\usepackage{xcolor}
\usepackage{fancyhdr}

% lastpage no es necesario para compilar, pero permite mostrar "pagina de N".
\newif\ifHayUltimaPagina
\IfFileExists{lastpage.sty}{
  \usepackage{lastpage}
  \HayUltimaPaginatrue
}{
  \HayUltimaPaginafalse
}

% -----------------------------------------------------------------------------
% CONFIGURACION DEL TRABAJO: editar solamente este bloque para reutilizarlo.
% -----------------------------------------------------------------------------
\newcommand{\NombreMateria}{Redes de Computadoras}
\newcommand{\TipoTrabajo}{Trabajo Práctico}
\newcommand{\NumeroTrabajo}{4}
% Cambiar por "Teórico" o "Práctico" según corresponda.
\newcommand{\TipoTP}{Teórico}
\newcommand{\NombreTrabajo}{Protocolos de Capa de Transporte}
\newcommand{\NombreGrupo}{NetRunners}
\newcommand{\Docente}{FACUNDO NICOLAS OLIVA CUNEO}
\newcommand{\LogoArchivo}{logo.png}
\newcommand{\LogoRuta}{\LogoArchivo}

% Una linea por integrante. Se pueden agregar o quitar filas.
\newcommand{\Integrantes}{%
  \begin{tabular}{ll}
    Bastida, Lucas Ramiro & DNI 39444511 \\
    Contessi, Ruben Andres & DNI 33315235 \\
    Garay, Ignacio Hernan & DNI 44191925 \\
    Giraudo, Juan Pablo & DNI 34970089 \\
    Peñaloza, Gonzalo Adrian & DNI 40441355 \\
    Quinteros del Castillo, Tomás Agustín & DNI 46169739 \\
    Vasconsellos Blason, Benjamin Alejandro & DNI 45642879 \\
  \end{tabular}%
}

% -----------------------------------------------------------------------------
% Elementos comunes de caratula y encabezado.
% -----------------------------------------------------------------------------
\newcommand{\InsertarLogo}[1][0.24\textwidth]{%
  \ifx\LogoRuta\empty
    \fbox{\parbox[c][1cm][c]{#1}{\centering Logo institucional\\\small (colocar archivo)}}%
  \else
    \includegraphics[width=#1,keepaspectratio]{\LogoRuta}%
  \fi
}

\newcommand{\TituloCorto}{TP \NumeroTrabajo (\TipoTP): \NombreTrabajo}

\pagestyle{fancy}
\fancyhf{}
\fancyhead[L]{\InsertarLogo[1.4cm]}
\fancyhead[C]{\small\textbf{\TituloCorto}}
\fancyhead[R]{\small\NombreGrupo}
\renewcommand{\headrulewidth}{0.4pt}
\fancyfoot[C]{\small Pagina \thepage\ifHayUltimaPagina\ de \pageref{LastPage}\fi}
\renewcommand{\footrulewidth}{0pt}

\begin{document}

\begin{titlepage}
  \thispagestyle{empty}
  \begin{center}
    \InsertarLogo[0.32\textwidth]

    \vfill
    {\Large\textsc{\NombreMateria}\par}
    \vspace{1.5cm}
    {\Huge\bfseries \TipoTrabajo\par}
    \vspace{0.35cm}
    {\LARGE N.\textsuperscript{o} \NumeroTrabajo (\TipoTP)\par}
    \vspace{0.8cm}
    {\Large\NombreTrabajo\par}

    \vfill
    \begin{tabular}{rl}
      \textbf{Grupo:} & \NombreGrupo \\
      \textbf{Docente:} & \Docente
    \end{tabular}

    \vspace{1cm}
    \begin{center}
      \textbf{Integrantes}\\[0.2cm]
      \Integrantes
    \end{center}

    \vfill
    {\large \FechaEnEspanol\par}
  \end{center}
\end{titlepage}

\begin{enumerate}[label=20.\arabic*, itemsep=0.3cm]

  \item 
  \item 
  \item 
  \textbf{Explique el uso de la multiplexación en el contexto de un protocolo de transporte.}

    La \textbf{multiplexación} permite que el protocolo de transporte utilice de manera más eficiente las conexiones de red. La idea es que varias conexiones de transporte puedan compartir una misma conexión de red, o que una conexión de transporte pueda utilizar varias conexiones de red.

    Se pueden distinguir dos casos:

    \begin{itemize}
    \item \textbf{Multiplexación hacia arriba (upward multiplexing):} varias conexiones de transporte utilizan una única conexión de red. Esto permite reducir la cantidad de conexiones de red necesarias, compartiendo una misma conexión entre diferentes conexiones de transporte.

    ```
    \item \textbf{Multiplexación hacia abajo (downward multiplexing):} una única conexión de transporte utiliza varias conexiones de red. Los datos pueden distribuirse entre estas conexiones, lo que permite aprovecharlas simultáneamente y puede mejorar el rendimiento.
    ```

    \end{itemize}

    En resumen la multiplexación permite aprovechar mejor los recursos de la red, ya sea haciendo que varias conexiones de transporte compartan una misma conexión de red, o utilizando varias conexiones de red para una misma conexión de transporte.
  \item 

    \textbf{Describa brevemente el esquema de créditos utilizado por TCP para el control de flujo.}

    TCP utiliza un \textbf{esquema de créditos} para realizar el control de flujo y evitar que el emisor envíe más datos de los que el receptor puede almacenar.

    Cada octeto de datos posee un número de secuencia, y los segmentos TCP incluyen un número de secuencia, un número de confirmación (\textit{ACK}) y un tamaño de ventana. El receptor utiliza la ventana para indicar al emisor cuántos octetos adicionales está preparado para recibir.

    Por ejemplo, si un ACK contiene \(AN=i\) y una ventana \(W=j\), significa que todos los octetos anteriores a \(i\) fueron recibidos correctamente, que se espera recibir el octeto \(i\) y que el emisor tiene permiso para enviar hasta \(j\) octetos adicionales.

    De esta manera, el receptor puede informar dinámicamente cuánto espacio disponible tiene en su buffer. Si la ventana es pequeña, el emisor debe reducir la cantidad de datos que envía; si aumenta, puede enviar más datos.

    Una característica importante del esquema de créditos de TCP es que \textbf{el control de flujo está desacoplado de las confirmaciones}. Es decir, un segmento puede confirmar datos recibidos sin necesariamente otorgar nuevos créditos, o puede otorgar créditos sin confirmar nuevos datos.

  \item 
  \item 
  \item 
  \item 
  \item 
  \item 
  \item 

\end{enumerate}

\end{document}
